\scanpage{91}
\mspage{8}
% archive sheet number: 90
\deleted{\unsure{1.25}}

1.26. \uline{Restriction aux revêtements modérément ramifiés} de $S$.

Comme on verra par la suite, le calcul des $\underline H^i_{\unsure{X_{\alpha0}|Y_\alpha}}(\Lambda_{\nu X_\alpha})$ peut se faire, parfois, à condition de se borner aux $S_\alpha$ qui sont modérément ramifiés sur $S$, correspondant donc aux \deleted{les} sous-extensions \inserted{finies} de l'extension $\bar K_t\subset\bar K$, \deleted{\unsure{leur}} correspondant au $p$-sous-groupe de Sylow $P$ de $I$. \deleted{\unsure{Puisque}} \deleted{Dans ce cas on voit} \deleted{La suite} La suite exacte de \unsure{Hochschild-Serre} pour $\deleted{X}U_{\bar K}\longrightarrow\deleted{X}U_{\bar K_t}$ \inserted{de gr $P$} donne alors
\[
H^{\unsure{*|\alpha}}(U_{\bar K_t},\Lambda_\nu)\Longleftarrow E_2^{p,q}=H^p(P,H^q(U_{\bar K},\Lambda_\nu)),
\]
et on a $E_2^{pq}=0$ si $p\ne0$, car $H^q(U_{\bar K},\Lambda_\nu)$ est de $\ell$-torsion, $\ell\ne p$. Donc on trouve
\[
H^*(U_{\bar K_t},\Lambda_\nu)=H^*(U_{\bar K},\Lambda_\nu)^P,
\tag*{(1.27)}
\]
invariants sous le groupe de $p$-Sylow. D'ailleurs, le premier \unsure{est aussi|est encore} isomorphe à $\varinjlim_\alpha H^*(U_{K_\alpha},\Lambda_\nu)$, où $K_\alpha$ parcourt les sous-extensions finies \uline{de} $\bar K_t$, \deleted{de} \unsure{ou|à} si on veut, les
\[
K(n)=K({\unsure{\pi|\Pi}}^{1/n})\qquad\text{pour }(p,n)=1.
\tag*{(1.28)}
\]
